\scanpage{38}
18) Les foncteurs entre cat. \unsure{triangulés|triangulées} \deleted{\unsure{doivent}} commutent au foncteur translation et transforment triangles \unsure{distingués} en \unsure{idem} ....
\[
D(C),\quad D^-(C)\qquad\text{importants }=\text{ sans qu'il y ait des projectifs.}
\]
\[
D^b(C)
\]
\[
Lf^*:D(Y)\longrightarrow D(X)
\]
quand on a $f:X\mathbin{-}Y$ morphisme \uuline{plat} de topos \uline{annelés}
\[
D^+(Y)\mathbin{-}D^+(X)
\]
\[
D^-(Y)\mathbin{-}D^-(X)
\]

$D^+(Y)\mathbin{-}D^+(X)$\qquad définition directe via résolutions injectives

\uline{ou} définition immédiate par ``passage au quotient''

[Pour calculer $Lf^*(K)$, pas la peine de remplacer au préalable $K$ par une résolution injective]

\uuline{NB}. Si $f$ n'est \uline{pas plat}, on \deleted{\unsure{y}} ne peut en général définir raisonnablement
\[
f^*:D^+(Y)\longrightarrow D^+(X)
\]
(de sorte à avoir des bonnes propriétés). Mais dans la suite de ces notes, on explique comment construire raisonnablement $Lf^*:D^-(Y)\mathbin{-}D^-(X)$, \inserted{``dérivé gauche'' de $f^*$,} en utilisant (non pas des rés. projectives mais) des résolutions plates. Ce n'est que si $f$ est à ``Tor-dim. finie'' que \unsure{ça} définit
\[
D^+(Y)\longrightarrow D^+(X)\ ;\quad D^b(Y)\mathbin{-}D^b(X).
\]
\margin{Rien n'empêche de construire un dérivé droit, $Rf^*$, mais il n'a pas le moindre intérêt.}
\mspage{37}
